% ============================================================================
%  Chapitre 3, partie IV : Validation de la dynamique
% ============================================================================

% ----------------------------------------------------------------------------
\section{Validation de la dynamique}
% ----------------------------------------------------------------------------

Cette partie sert à vérifier la dynamique obtenue par les équations
d'Euler-Lagrange-d'Alembert à l'aide des équations de Gibbs-Appell.
Cependant, celles obtenues par Gibbs-Appell ne font pas apparaître explicitement les forces de réaction, mais on peut les retrouver grâce aux relations cinématiques déjà
imposées. Si les deux approches donnent le même résultat, la dynamique
est validée et on peut commencer à travailler avec.

% ----------------------------------------------------------------------------
\subsection{Loi de Newton en \texorpdfstring{$x$ et $y$}{x et y}}

Dans le plan horizontal, la seule force appliquée à la sphère est le
frottement exercé par le sol en $P$, de composantes $(N_x,N_y)$. La
loi de Newton en $x$ et $y$ donne
\[
    m\,\ddot x = N_x,
    \qquad
    m\,\ddot y = N_y
\]
D'après le roulement, $\dot r = -\,\mat{A}_\alpha\,\dot\alpha$, donc
\[
    \begin{pmatrix} N_x\\ N_y \end{pmatrix}
    = m\,\ddot r
    = -\,m\,\bigl(\mat{A}_\alpha\,\ddot\alpha
    + \dot{\mat{A}}_\alpha\,\dot\alpha\bigr)
\]
Les frottements sont donc connus dès que l'on connaît $\alpha(t)$. On retrouve
$\vect{\lambda} = (N_x,N_y)$ de II.3. En remplaçant $m\,\ddot r$ par
$\vect{\lambda}$ dans les équations de Gibbs-Appell, et en y ajoutant
la loi de Newton en $x$ et $y$, on retrouve exactement les équations
d'Euler-Lagrange-d'Alembert.

% ----------------------------------------------------------------------------
\subsection{Résumé : Dynamique du système}
\begin{theoreme}[Dynamique de la sphère de Chaplygin]
    Avec
    \[
        \left\{
        \begin{aligned}
            \omega_1 &= \dot\varphi - \dot\psi\,s_\theta\\
            \omega_2 &= \dot\theta\,c_\varphi
                + \dot\psi\,s_\varphi c_\theta\\
            \omega_3 &= -\dot\theta\,s_\varphi
                + \dot\psi\,c_\varphi c_\theta
        \end{aligned}
        \right.
        \qquad
        \left\{
        \begin{aligned}
            \Gamma_1 &= I_1\,\dot\omega_1
                + (I_3 - I_2)\,\omega_2\,\omega_3\\
            \Gamma_2 &= I_2\,\dot\omega_2
                + (I_1 - I_3)\,\omega_3\,\omega_1\\
            \Gamma_3 &= I_3\,\dot\omega_3
                + (I_2 - I_1)\,\omega_1\,\omega_2
        \end{aligned}
        \right.
    \]
    l'orientation de la sphère vérifie
    \begin{align*}
        &\Gamma_1 + m\,a\,c_\theta\,(s_\psi\,\ddot x - c_\psi\,\ddot y)
            = 0\\
        &c_\varphi\,\Gamma_2 - s_\varphi\,\Gamma_3
            + m\,a\,(c_\psi\,\ddot x + s_\psi\,\ddot y) = 0\\
        &-\,s_\theta\,\Gamma_1 + s_\varphi c_\theta\,\Gamma_2
            + c_\varphi c_\theta\,\Gamma_3 = 0
    \end{align*}
    la position du centre vérifie le roulement
    \[
        \left\{
        \begin{aligned}
            \dot x &= a\,(\dot\varphi\,s_\psi c_\theta
                + \dot\theta\,c_\psi)\\
            \dot y &= -\,a\,(\dot\varphi\,c_\psi c_\theta
                - \dot\theta\,s_\psi)
        \end{aligned}
        \right.
    \]
    et les frottements sont donnés par la loi de Newton
    \[
        N_x = m\,\ddot x,
        \qquad
        N_y = m\,\ddot y
    \]
\end{theoreme}

% ----------------------------------------------------------------------------
\subsection{Remarques}

\begin{remarque}[Complémentarité des deux méthodes]
    Les deux méthodes sont complémentaires et chacune a ses points
    forts, et on ne peut pas dire que l'une peut l'emporter sur l'autorisées.
    
    L'approche d'Euler-Lagrange-d'Alembert est directe : elle
    donne toutes les inconnues, frottements compris, mais à un prix de
    calculs lourds et longs. 
    
    L'approche de Gibbs-Appell donne plus rapidement
    les équations en $\alpha$, mais ne fait pas apparaître les
    réactions. Il suffit alors d'utiliser les relations cinématiques
    imposées pour retrouver la dynamique complète du système.
\end{remarque}

\begin{remarque}[Recul]
    Avec le recul du chapitre 4 (je fais la validation après avoir fait le modèle de simulation), l'approche de Gibbs-Appell m'aurait
    permis d'obtenir plus rapidement l'équation en $\ddot\alpha$. 
    
    En revanche, sans l'approche d'Euler-Lagrange-d'Alembert, je n'aurais
    pas vu qu'il était important de faire intervenir les frottements,
    dont dépend la validité de la condition de non-glissement. 
\end{remarque}

